\chapter{Metodologia}
\label{cap:metodologia}

\section{Caracterização da pesquisa}
\label{sec:caracterizacao}

Quanto à natureza, a pesquisa é aplicada, pois trata de um problema concreto: a organização do registro e do acompanhamento das informações do QRQC em uma rotina industrial. Quanto aos objetivos, é descritiva, uma vez que busca registrar como a aplicação estrutura esse processo e como os dados por ela gerados são apresentados, sem estabelecer relações de causa e efeito \cite{prodanov2013}.

Quanto à abordagem, a pesquisa combina elementos qualitativos e quantitativos. A parte qualitativa consiste no exame das páginas, do modelo de dados e das regras da aplicação, interpretados à luz das etapas do QRQC descritas na literatura. A parte quantitativa consiste na descrição dos indicadores do QRQC apresentados em uma ferramenta de \textit{Business Intelligence} (BI), alimentada pelos registros da própria aplicação \cite{prodanov2013}.

Quanto aos procedimentos, adota-se o estudo de caso único, adequado quando se pretende examinar um fenômeno contemporâneo em seu contexto real, com o apoio de múltiplas fontes de evidência \cite{yin2018}. O caso é o QRQC digitalizado por meio de uma aplicação desenvolvida em Mendix. A unidade de análise é essa aplicação, examinada em sua totalidade quanto ao suporte que oferece às etapas do QRQC; funcionalidades de apoio, como cadastros de usuários e de tipos de defeito, são consideradas na medida em que sustentam o registro e o acompanhamento das ocorrências.

O estudo não compara o processo antes e depois da implantação da aplicação, pois não há registros suficientes sobre a rotina anterior. Assim, os indicadores obtidos no BI descrevem o funcionamento do processo no período analisado, mas não permitem afirmar que a digitalização tenha provocado melhorias de desempenho. A organização é identificada por nome fictício, e os dados apresentados são anonimizados.
